\section{Reversing symmetry and the augmentation determinant}

Let $p>3$ satisfy $p\equiv1\pmod3$, and fix
$\rho\in\Fp^\times$ of order three.  On
$\Hh_p=\ell^2(\Fp)$ define
\[
(U_pf)(Q)=p^{-1/2}\sum_{q\in\Fp}
\exp\!\left(\frac{2\pi i}{p}(qQ+2q^3)\right)f(q),
\qquad
(R_pf)(x)=f(\rho x).
\]
The operator $U_p$ is unitary because it is a normalized finite Fourier
transform followed by a unitary chirp.

\begin{theorem}[Intrinsic cubic reversing symmetry]\label{thm:symmetry}
Let $g(q,p_0)=(\rho q,\rho^{-1}p_0)$.  Then
\[
H_0g=g^{-1}H_0,
\qquad
U_pR_p=R_p^{-1}U_p.
\]
Consequently $T_p=U_p^2$ commutes with $R_p$.
\end{theorem}

\begin{proof}
The classical identity follows by direct substitution.  For the quantum
identity, change the summation variable by $q=\rho^{-1}r$ and use
$\rho^3=1$.  Squaring the relation gives $[U_p^2,R_p]=0$.
\end{proof}

Let $\omega=e^{2\pi i/3}$ and
$\Hh_{p,k}=\ker(R_p-\omega^kI)$.  Multiplication by $\rho$ has the fixed
orbit $\{0\}$ and $(p-1)/3$ free orbits, giving
\[
d_{p,0}=\frac{p+2}{3},\qquad
d_{p,1}=d_{p,2}=\frac{p-1}{3}.
\]
Moreover, $U_p$ maps grade $k$ to grade $-k$ and intertwines the
restrictions of $T_p$.  Thus $T_{p,1}$ and $T_{p,2}$ are unitarily
equivalent.

\begin{definition}[Integral augmentation factor]
For $D_{p,k}(z)=\det(I-zT_{p,k})$, set
\[
\Daug_p(z)=
\frac{D_{p,0}(z)^2}{D_{p,1}(z)D_{p,2}(z)}
=\left(\frac{D_{p,0}(z)}{D_{p,1}(z)}\right)^2.
\]
\end{definition}

The weights $(2,-1,-1)$ lie in the integral representation ring.  This is an
equivariant virtual determinant, not a Berezinian: a $\mathbb Z/3$ grading
does not provide a canonical $\mathbb Z/2$ parity.  Its virtual dimension is
two, but the individual degrees are of order $p$.

\begin{proposition}[Determinant normalization]\label{prop:detnorm}
For $k=0,1,2$,
\[
\det T_{p,k}=(-1)^{(p-1)/6}.
\]
Consequently $D_{p,0}/D_{p,1}\sim-z$ and
$\Daug_p(z)\sim z^2$ as $z\to\infty$.
\end{proposition}

The proof is given in \cref{app:detnorm}.  In particular, no cancellation can
make the augmentation factor constant.
